\chapter[Behavioral Sources in LTspice]
{Mathematical Modelling with Behavioral Sources in LTspice}
\label{app:behavioral-sources}
\index{LTspice!behavioral source}
\index{behavioral source}
\index{controlled source}

\section{Behavioral E and G Sources}
\label{sec:behavioral-eg-sources}

Behavioral voltage and current sources make it possible to implement
mathematical models directly in LTspice. Node voltages, branch currents,
parameters, arithmetic operations, logical operations, and built-in functions
can be combined in one expression.

\begin{formulaBox}{Basic LTspice syntax}
An E source generates a voltage between its output nodes:
\begin{lstlisting}[style=LTspice]
E1 node1 node2 VALUE={mathematical_expression}
\end{lstlisting}

A G source generates a branch current between its output nodes:
\begin{lstlisting}[style=LTspice]
G1 node1 node2 VALUE={mathematical_expression}
\end{lstlisting}

Typical expression operands are
\begin{lstlisting}[style=LTspice]
V(node)              ; node voltage relative to ground
V(node_plus,node_minus) ; differential node voltage
I(Vsense)             ; current through a voltage source
Ix(U1:pin)            ; current into a subcircuit pin
\end{lstlisting}
\end{formulaBox}

\begin{noteBox}
\textbf{Elements used in the toolboxes.}
The principal element designators are
\texttt{V}, \texttt{I}, \texttt{L}, \texttt{C}, \texttt{R}, \texttt{K},
\texttt{E}, \texttt{G}, \texttt{S}, \texttt{T}, and \texttt{O}.
Consult the LTspice help for the precise syntax, sign convention, and analysis
restrictions of each element.
\end{noteBox}

\begin{important}{Dimensional consistency}{imp:behavioral-source-units}
The value returned by an E-source expression is interpreted as volts; the
value returned by a G-source expression is interpreted as amperes. LTspice
does not enforce physical dimensional consistency inside a behavioral
expression. The model author must therefore verify all scaling factors and
units.
\end{important}

\section{Functions in Behavioral Expressions}
\label{sec:behavioral-functions}

The following functions can be used in behavioral E and G sources. Optional
arguments are shown in square brackets. Function availability can depend on the
installed LTspice version.

\begingroup
\small
\setlength{\LTpre}{0pt}
\setlength{\LTpost}{6pt}
\renewcommand{\arraystretch}{1.12}
\begin{longtable}{>{\ttfamily\raggedright\arraybackslash}p{0.31\textwidth}
                  >{\raggedright\arraybackslash}p{0.61\textwidth}}
\caption{Functions available in LTspice behavioral expressions.}
\label{tab:ltspice-behavioral-functions}\\
\toprule
\normalfont\bfseries Function & \bfseries Description \\
\midrule
\endfirsthead
\caption[]{Functions available in LTspice behavioral expressions (continued).}\\
\toprule
\normalfont\bfseries Function & \bfseries Description \\
\midrule
\endhead
\midrule
\multicolumn{2}{r}{\small Continued on next page}\\
\endfoot
\bottomrule
\endlastfoot
abs(x) & Absolute value of $x$. \\
absdelay(x,t[,tmax]) & Delays $x$ by $t$; the optional maximum delay is \texttt{tmax}. \\
acos(x) & Real part of the arc cosine. For example, \texttt{acos(-5)} returns \texttt{3.14159}, not a complex result. \\
arccos(x) & Synonym for \texttt{acos(x)}. \\
acosh(x) & Real part of the inverse hyperbolic cosine. \\
asin(x) & Real part of the arc sine. Out-of-domain arguments are limited to the real result. \\
arcsin(x) & Synonym for \texttt{asin(x)}. \\
asinh(x) & Inverse hyperbolic sine. \\
atan(x) & Arc tangent of $x$. \\
arctan(x) & Synonym for \texttt{atan(x)}. \\
atan2(y,x) & Four-quadrant arc tangent of $y/x$. \\
atanh(x) & Inverse hyperbolic tangent. \\
buf(x) & Returns 1 if $x>0.5$; otherwise returns 0. \\
cbrt(x) & Real cube root of $x$. \\
ceil(x) & Smallest integer greater than or equal to $x$. \\
cos(x) & Cosine of $x$; the argument is in radians. \\
cosh(x) & Hyperbolic cosine of $x$. \\
ddt(x) & Time derivative $\diff x/\diff t$. \\
delay(x,t[,tmax]) & Same as \texttt{absdelay(x,t[,tmax])}. \\
dnlim(x,A,Delta) & Smooth approximation of \texttt{max(x,A)} with first-derivative transition width $\Delta$. \\
exp(x) & Exponential function $\mathrm{e}^{x}$. \\
floor(x) & Largest integer less than or equal to $x$. \\
hypot(x,y) & Returns $\sqrt{x^2+y^2}$ with suitable numerical scaling. \\
idt(x[,ic[,a]]) & Integrates $x$; \texttt{ic} is the optional initial condition and the integrator is reset when \texttt{a} is true. \\
idtmod(x[,ic[,m[,o]]]) & Integrates $x$, resets on reaching modulus \texttt{m}, and offsets the output by \texttt{o}. \\
if(x,y,z) & Returns $y$ if $x>0.5$; otherwise returns $z$. \\
int(x) & Converts $x$ to an integer. \\
inv(x) & Returns 0 if $x>0.5$; otherwise returns 1. \\
limit(x,y,z) & Returns the intermediate value of $x$, $y$, and $z$. \\
ln(x) & Natural logarithm of $x$. \\
log(x) & Alternate syntax for \texttt{ln(x)}. \\
log10(x) & Base-10 logarithm of $x$. \\
max(x,y) & Greater of $x$ and $y$. \\
min(x,y) & Smaller of $x$ and $y$. \\
mod(x,y) & Remainder of the division $x/y$. \\
noise() & Returns 1 during a \texttt{.noise} analysis; otherwise returns 0. \\
pow(x,y) & Real part of $x^y$; for example, \texttt{pow(-1,0.5)} returns 0 rather than an imaginary value. \\
pwr(x,y) & Returns $|x|^y$. \\
pwrs(x,y) & Returns $\operatorname{sgn}(x)|x|^y$. \\
rand(x) & Pseudorandom value between 0 and 1 determined by the integer value of $x$. \\
random(x) & Similar to \texttt{rand(x)}, but transitions smoothly between values. \\
round(x) & Integer nearest to $x$. \\
sdt(x[,ic[,assert]]) & Alternate syntax for \texttt{idt(x[,ic[,assert]])}. \\
sgn(x) & Sign of $x$: $-1$, 0, or 1. \\
sin(x) & Sine of $x$; the argument is in radians. \\
sinh(x) & Hyperbolic sine of $x$. \\
smallsig() & Returns 1 during a \texttt{.ac} or \texttt{.noise} analysis; otherwise returns 0. \\
square(x) & Returns $x^2$. \\
sqrt(x) & Square root of $x$. \\
table(x,a,b,c,d,...) & Linearly interpolates a value for $x$ from the point pairs $(a,b),(c,d),\ldots$. \\
table(x, .include ``filename'') & Interpolates from point pairs stored as numbers in an included file. \\
tan(x) & Tangent of $x$. \\
tanh(x) & Hyperbolic tangent of $x$. \\
u(x) & Unit step: 1 if $x>0$; otherwise 0. \\
uplim(x,A,Delta) & Smooth approximation of \texttt{min(x,A)} with first-derivative transition width $\Delta$. \\
uramp(x) & Returns $x$ if $x>0$; otherwise returns 0. \\
white(x) & Pseudorandom value between $-0.5$ and $0.5$ with smoother transitions than \texttt{random(x)}. \\
xor(x,y) & Converts $x$ and $y$ to Boolean values and returns their exclusive OR. \\
!(x) & Alternate syntax for \texttt{inv(x)}. \\
\textasciitilde(x) & Alternate syntax for \texttt{inv(x)}. \\
\end{longtable}
\endgroup

\section{Operators and Precedence}
\label{sec:behavioral-operators}

Table~\ref{tab:ltspice-operators} lists the behavioral-expression operators in
ascending order of precedence. Operators in the same group have equal
precedence; parentheses should be used whenever the intended evaluation order
could be ambiguous.

\begingroup
\small
\renewcommand{\arraystretch}{1.14}
\begin{longtable}{>{\ttfamily\raggedright\arraybackslash}p{0.19\textwidth}
                  >{\raggedright\arraybackslash}p{0.73\textwidth}}
\caption{LTspice behavioral-expression operators in ascending order of precedence.}
\label{tab:ltspice-operators}\\
\toprule
\normalfont\bfseries Operator & \bfseries Description \\
\midrule
\endfirsthead
\caption[]{LTspice behavioral-expression operators (continued).}\\
\toprule
\normalfont\bfseries Operator & \bfseries Description \\
\midrule
\endhead
\bottomrule
\endfoot
\& & Converts both operands to Boolean values and returns logical AND. \\
| & Converts both operands to Boolean values and returns logical OR. \\
\textasciicircum & Converts both operands to Boolean values and returns logical XOR. \\
>, < & Greater-than and less-than comparisons. \\
>=, <= & Greater-than-or-equal and less-than-or-equal comparisons. \\
== & Returns true if both expressions are equal. \\
!=, <> & Returns true if both expressions are unequal. \\
+, - & Floating-point addition and subtraction. \\
*, / & Floating-point multiplication and division. \\
\% & Floating-point remainder; equivalent to \texttt{mod(x,y)}. \\
** & Raises the left operand to the power of the right operand. Only the real part is returned. \\
!, \textasciitilde & Unary Boolean inversion of the following expression. \\
\end{longtable}
\endgroup

\begin{noteBox}
\textbf{Boolean conversion.}
For the behavioral functions \texttt{buf()}, \texttt{if()}, and
\texttt{inv()}, LTspice treats values greater than 0.5 as true. The exact
threshold behavior of logical operators should be checked when expressions can
take non-binary analog values.
\end{noteBox}

\section{Example: Single-Phase Transformer}
\label{sec:behavioral-transformer-example}

The example implements the voltage equations by numerical integration and
solves the flux-current relations as a simultaneous linear system.

\subsection{Differential equations}

The winding-voltage equations are
\begin{align}
u &= R_1 i_1+\frac{\diff\psi_1}{\diff t},\\
0 &= (R_2+R_{\mathrm{load}})i_2
     +\frac{\diff\psi_2}{\diff t}.
\end{align}

Consequently,
\begin{align}
\frac{\diff\psi_1}{\diff t}&=u-R_1i_1,\\
\frac{\diff\psi_2}{\diff t}&=-(R_2+R_{\mathrm{load}})i_2.
\end{align}

\subsection{Flux-current equations}

The coupled linear system is
\begin{equation}
\begin{bmatrix}\psi_1\\\psi_2\end{bmatrix}
=L_{1h}
\begin{bmatrix}
1+\sigma_1 & 1\\
1 & 1+\sigma_2
\end{bmatrix}
\begin{bmatrix}i_1\\i_2\end{bmatrix}.
\label{eq:behavioral-transformer-matrix}
\end{equation}

Equivalently, the two residual equations are
\begin{align}
0&=\psi_1-(1+\sigma_1)L_{1h}i_1-L_{1h}i_2,\\
0&=\psi_2-L_{1h}i_1-(1+\sigma_2)L_{1h}i_2.
\end{align}

\subsection{LTspice netlist}

\begin{lstlisting}[style=LTspice,caption={Behavioral implementation of a single-phase transformer.},label={lst:behavioral-transformer}]
* Single-phase transformer using behavioral sources
* Parameters
.PARAM U=20 R1=10 R2=8 R_Last=17
.PARAM L1h=0.1 sigma1=0.1 sigma2=0.15

* Integrator for flux psi1.
* G-C-R realization is used because idt() can be less robust.
Gpsi1  0 psi1 VALUE={ U - R1*V(i1) }
Cpsi1  psi1 0 1 IC=0
Rpsi1  psi1 0 1E6

* Integrator for flux psi2
Gpsi2  0 psi2 VALUE={ -(R2+R_Last)*V(i2) }
Cpsi2  psi2 0 1 IC=0
Rpsi2  psi2 0 1E6

* Simultaneous linear system
Gi1  0 i1 VALUE={ V(psi1)
+ (-(1+sigma1)*L1h*V(i1)-L1h*V(i2)) }
Ri1  i1 0 1E12

Gi2  0 i2 VALUE={ V(psi2)
+ (-L1h*V(i1)-(1+sigma2)*L1h*V(i2)) }
Ri2  i2 0 1E12

.TRAN 0.001 0.04 0 0.001 UIC
\end{lstlisting}

\begin{remark}{Solution of the algebraic system}{rem:behavioral-linear-system}
The resistances \texttt{Ri1} and \texttt{Ri2} are set to
\SI{1e12}{\ohm}. The G-source currents can therefore produce only negligible
node voltages unless their behavioral expressions are nearly zero. During each
Newton iteration, LTspice consequently adjusts \texttt{V(i1)} and
\texttt{V(i2)} until both residual equations are satisfied simultaneously.
The node voltages thus represent the calculated currents $i_1$ and $i_2$.
\end{remark}

\begin{important}{Numerical interpretation of the G-C-R integrator}{imp:gcr-integrator}
With \texttt{Cpsi1=1 F}, the capacitor relation gives
$\diff V(\texttt{psi1})/\diff t=I(\texttt{Gpsi1})$, so the node voltage
directly represents flux linkage. The parallel \SI{1}{\mega\ohm} resistor
provides a DC path but also introduces a small leakage term. Increase its value
if this artificial decay is significant, while retaining a sufficiently robust
DC operating point.
\end{important}

\begin{chaptersummary}[Appendix Summary]
Behavioral E and G sources allow differential equations, algebraic equations,
logical functions, and lookup tables to be implemented directly in LTspice.
E sources return voltages, whereas G sources return currents. Integrators may
be formed either with \texttt{idt()} or explicitly with a G-C-R network. A
simultaneous algebraic system can be solved by expressing each residual with a
behavioral G source and forcing it toward zero through a very large shunt
resistance. Careful attention to sign conventions, units, damping, and
operator precedence is essential for numerically robust models.
\end{chaptersummary}
